% Problem-section file following ../_template.tex.
\section{Parallel versus nested-adaptive superchannel discrimination}
\label{sec:problem-66}

\paragraph{Problem.}
Can nested-adaptive strategies improve the Stein exponent for discriminating
two finite-dimensional quantum superchannels?  A superchannel maps channels
$\mathcal N:A\to B$ to channels $C\to D$ and admits a realization
\begin{equation}
  \Theta(\mathcal N)
  =\mathcal D\circ(\mathcal N\otimes\operatorname{id}_S)\circ\mathcal E,
  \qquad
  \mathcal E:C\to A S,
  \quad \mathcal D:B S\to D.
  \label{eq:p66-superchannel-realization}
\end{equation}
The memory system $S$ in Eq.~\eqref{eq:p66-superchannel-realization} is
internal to the superchannel.  In a fully parallel $n$-use strategy, one applies
$\Theta_i^{\otimes n}$ to a joint $n$-partite inserted channel and then tests
the resulting output on a joint input state.  A nested-adaptive strategy may
instead recursively insert the channel produced by one tested use into the
channel slot of another, with arbitrary compatible CPTP maps between uses and
a final binary measurement.  For
$\mathsf S\in\{\mathrm{par},\mathrm{nest}\}$, let
\begin{equation}
  \begin{aligned}
  \beta_{\varepsilon,n}^{\mathsf S}(\Theta_1\|\Theta_2)
    &:=\inf\{\beta_n(P):P\in\mathsf S_n,\ \alpha_n(P)\leq\varepsilon\},\\
  \zeta_{\mathsf S}(\Theta_1\|\Theta_2)
    &:=\lim_{\varepsilon\downarrow0}\liminf_{n\to\infty}
      -\frac1n\log_2
      \beta_{\varepsilon,n}^{\mathsf S}(\Theta_1\|\Theta_2),
  \end{aligned}
  \label{eq:p66-stein-exponents}
\end{equation}
where $\alpha_n$ and $\beta_n$ are the type-I and type-II errors.  Is
\begin{equation}
  \zeta_{\mathrm{nest}}(\Theta_1\|\Theta_2)
  =\zeta_{\mathrm{par}}(\Theta_1\|\Theta_2)
  \label{eq:p66-nested-parallel-equality}
\end{equation}
for every pair $\Theta_1,\Theta_2$?

\subsection*{Status}

Unsolved

\subsection*{Source}

Hirche explicitly conjectured the equality in
Eq.~\eqref{eq:p66-nested-parallel-equality} while developing the first
asymptotic discrimination framework for quantum superchannels
\sourcecite{ref:p66-hirche}{Hir23}.

\subsection*{Progress}

\begin{itemize}
  \item The fully parallel Stein exponent in
  Eq.~\eqref{eq:p66-stein-exponents} is known exactly:
  \begin{equation}
    \zeta_{\mathrm{par}}(\Theta_1\|\Theta_2)
    =D_{\rm sc}^{\infty}(\Theta_1\|\Theta_2),
    \label{eq:p66-parallel-exponent}
  \end{equation}
  In Eq.~\eqref{eq:p66-parallel-exponent}, $D_{\rm sc}^{\infty}$ is the
  regularized superchannel relative
  entropy optimized over joint inserted channels and input states
  \sourcecite{ref:p66-hirche}{Hir23}.

  \item Every parallel strategy can be embedded into a nested one.  Hirche's
  amortized-superchannel meta-converse therefore gives
  \begin{equation}
    D_{\rm sc}^{\infty}(\Theta_1\|\Theta_2)
    \leq\zeta_{\mathrm{nest}}(\Theta_1\|\Theta_2)
    \leq D_{\rm sc}^{A}(\Theta_1\|\Theta_2),
    \label{eq:p66-known-hierarchy}
  \end{equation}
  with $D_{\rm sc}^{A}$ the amortized superchannel relative entropy
  \sourcecite{ref:p66-hirche}{Hir23}.

  \item The channel chain rule underlying
  Eq.~\eqref{eq:p66-known-hierarchy} implies
  $D_{\rm sc}^{A}=D_{\rm sc}^{\infty}$, and hence
  Eq.~\eqref{eq:p66-nested-parallel-equality}, conditional on the
  diamond-smoothed channel AEP in Problem~\ref{sec:problem-65}.  No
  unconditional proof or counterexample is known
  \sourcecite{ref:p66-hirche}{Hir23}.

  \item For classical superchannels, all relevant relative-entropy
  amortizations collapse and product strategies already achieve the optimum;
  thus Eq.~\eqref{eq:p66-nested-parallel-equality} holds in that special case
  \sourcecite{ref:p66-hirche}{Hir23}.
\end{itemize}

\subsection*{References}

\begin{enumerate}[label={},leftmargin=4.8em,labelsep=0.5em]
  \footnotesize
  \item[\textup{[Hir23]}]\label{ref:p66-hirche}
  C. Hirche, ``Quantum Network Discrimination,''
  \emph{Quantum} \textbf{7}, 1064 (2023).
  \href{https://doi.org/10.22331/q-2023-07-25-1064}%
       {doi:10.22331/q-2023-07-25-1064};
  \href{https://arxiv.org/abs/2103.02404}{arXiv:2103.02404}.
\end{enumerate}

\subsection*{Comment}

This problem compares two causal strategy classes.  It is distinct from
Problem~\ref{sec:problem-67}, which allows arbitrary interleavings of the
physical components of different superchannel uses.

\subsection*{Tag}

Channel discrimination; Quantum hypothesis testing; Quantum networks;
Superchannels

\subsection*{ID}

\texttt{op\_a4600b38b94042a8}
